% ============================================================================
% Chapter 17 --- How to read this playbook
% ----------------------------------------------------------------------------
% Purpose : establish vocabulary, dependencies, trajectory selection and
%           pace options for the migration playbook. Opens Part III.
% Page target: ~4 pages.
% Dependencies: Ch. 4 (self-assessment feeds the trajectory selection),
%               Ch. 16 (conformance suite closes each wave),
%               Ch. 18--27 (the wave chapters apply the framework here).
% ----------------------------------------------------------------------------
% Decisions registered in _coherence-EN.md:
%   D17.1 stable vocabulary for Part III (frozen)
%   D17.2 ten waves: Wave 0 to Wave 9
%   D17.3 dependency map between waves
%   D17.4 trajectory selection matrix consistent with D4.5 (five profiles)
%   D17.5 three pace options: 5y / 7y / 10y
%   D17.6 reversibility philosophy: every wave reversible at acceptable
%         cost during a published window
% ============================================================================

\chapter{How to read this playbook}
\label{ch:playbook-reading}

\begin{caveatbox}[Dependencies]
\noindent The playbook of Part~III applies the architecture of
Part~II under the framework of Part~I. It expects the reader to
have completed the self-assessment of \trchap{ch:self-assessment}{},
and to be familiar with the sovereignty grid of
\trchap{ch:sovereignty-grid}{} and with the conformance suite of
\trchap{ch:conformance}{}, which closes every wave.
\end{caveatbox}

The playbook is structured as ten waves (Wave~0 to Wave~9) that, in
combination, take an institution from its current brownfield
position to the target architecture of Part~II. The waves are not
calendar phases: their sequencing is governed by their pre-requisite
graph, and their pace is a function of the institution's
operational rhythm rather than of an externally imposed timeline.

%-----------------------------------------------------------------------------
\section{Objectives of the playbook}
\label{sec:pr-objectives}

The playbook serves four institutional objectives that the wave
structure of Chapters~18--27 then operationalises.

\begin{enumerate}[leftmargin=2em,nosep]
  \item \textbf{Move from brownfield to target.} The principal
    objective is to take the institution from its current
    brownfield position --- whatever that may be --- to the
    target architecture of Part~II. The waves are the units of
    that movement, and the institution's chosen trajectory
    (\S\ref{sec:pr-trajectories}) determines how that movement
    is paced.
  \item \textbf{Operationalise the sovereignty grid.} The grid of
    \trchap{ch:sovereignty-grid}{} is a deliberation instrument;
    the playbook makes it operational by turning it into the
    gating mechanism at procurement, at wave closure, and at
    every renewal. A wave that closes against the grid produces
    a documented sovereignty profile per affected component; a
    procurement that follows the playbook is informed by the
    grid before it is decided.
  \item \textbf{Preserve reversibility.} Every wave is reversible
    at acceptable cost during a published window. The playbook
    treats reversibility as non-negotiable
    (\trchap{ch:reversibility}{}). The institution that follows
    the playbook does not commit to a one-way change.
  \item \textbf{Generate a longitudinal record.} Conformance
    results, post-mortems, and re-assessment outcomes accumulate
    into a record that the institution uses for its own
    governance and that the international peer-comparison phase
    of this initiative depends on. The playbook is consequently
    not only a transformation tool but also a documentation
    discipline.
\end{enumerate}

\paragraph{What the playbook does not promise.} The playbook does
not promise that every institution will reach the target
architecture in a fixed period, that every wave will succeed at
its first attempt, or that the architecture will be operationally
identical at every adopting institution. The first two are
managed through the pace options of \S\ref{sec:pr-pace} and the
reversibility playbook of \trchap{ch:reversibility}{}; the third
is a feature, since the sovereignty grid permits institutional
differentiation within the shared contract.

\paragraph{Failure mode to avoid.} The most common failure mode
of architecture-transformation programmes is the substitution of
products without architectural progress: a new identity provider
replaces an old one, but the institution has not gained
sovereignty, replaceability, or auditability. The playbook
guards against this through the conformance suite of
\trchap{ch:conformance}{}: a wave closes only when the
conformance subset passes, and conformance verifies that the
contracts of Part~II hold --- not merely that products have
changed.

%-----------------------------------------------------------------------------
\section{Vocabulary}
\label{sec:pr-vocabulary}

The chapters of Part~III use the following terms in stable senses.

\begin{description}[leftmargin=1.5em,style=nextline,nosep]
  \item[Wave.] One discrete stage of the migration, with a stated
    goal, declared pre-requisites, a coexistence pattern that
    governs the period of the wave, exit criteria that close it,
    and reversibility criteria that allow it to be undone. Waves
    are numbered Wave~0 to Wave~9.
  \item[Trajectory.] One of the five migration paths through the
    waves, distinguished by the institution's starting position:
    \emph{greenfield}, \emph{brownfield single-suite-centric},
    \emph{brownfield hybrid}, \emph{selective adoption}, or
    \emph{deferred}. Each wave chapter develops the three
    principal trajectories (greenfield, brownfield
    single-suite-centric, brownfield hybrid) in detail, with selective
    adoption and deferred treated as cases that diverge from the
    framework.
  \item[Coexistence pattern.] A documented bridging arrangement
    between an incumbent and an incoming component during a wave.
    Coexistence is the operational state of the wave; it is not a
    transitional inconvenience but the principal engineering effort
    of the wave's duration.
  \item[Pre-requisite.] An institutional or technical condition
    that must hold before the wave begins. Each wave declares its
    pre-requisites in the wave chapter; the dependency map of
    \S\ref{sec:pr-dependencies} consolidates them.
  \item[Exit criteria.] The conditions under which the wave is
    considered closed. Exit criteria include the relevant subset of
    the conformance suite of \trchap{ch:conformance}{}; they are
    not subjective.
  \item[Reversibility criteria.] The conditions under which the
    wave can be reversed at acceptable cost during a published
    window. Reversibility is non-negotiable in this playbook
    (\trchap{ch:reversibility}{}).
\end{description}

%-----------------------------------------------------------------------------
\section{Wave dependency map}
\label{sec:pr-dependencies}

The waves do not unfold in strict sequence; they unfold under a
dependency graph that allows certain waves to proceed in parallel
and others to wait. The reference dependency structure is given in
Table~\ref{tab:pr-deps}.

\begin{table}[H]
\centering\small
\caption{Wave dependency structure. ``Required'' pre-requisites must
be closed; ``Recommended'' pre-requisites strengthen the wave but do
not block it.}
\label{tab:pr-deps}
\begin{tabularx}{\textwidth}{%
  >{\hsize=0.55\hsize\linewidth=\hsize\bfseries\raggedright\arraybackslash}X%
  >{\hsize=1.50\hsize\linewidth=\hsize\raggedright\arraybackslash}X%
  >{\hsize=0.95\hsize\linewidth=\hsize\raggedright\arraybackslash}X%
}
\toprule
Wave & Required pre-requisites & Recommended \\
\midrule
0 — Foundations          & --- (entry point)                         & --- \\
1 — Identity             & 0                                         & --- \\
2 — Observability        & 0                                         & --- \\
3 — Network              & 0                                         & 2 (to baseline flows) \\
4 — Policy               & 1, 2                                      & 3 \\
5 — Pilot                & 1, 2, 3, 4                                & --- \\
6 — Endpoint             & 1, 4                                      & 2, 5 \\
7 — Audit \& IR          & 2, 4, 5                                   & --- \\
8 — Workload             & 1, 3, 4                                   & 5 \\
9 — Lifecycle            & 1--8 in some form                         & --- \\
\bottomrule
\end{tabularx}
\end{table}

\noindent Wave~0 is the entry point for every trajectory: it
establishes the governance and inventory without which the other
waves cannot be evaluated. Waves 1, 2 and 3 (Identity,
Observability, Network) can proceed in parallel after Wave~0 if the
institution's available staffing profile provides the operational
capacity to staff them. Wave~4
(Policy) requires both Identity and Observability. Wave~5 (Pilot)
requires the four preceding waves and is the first opportunity to
exercise the architecture end-to-end on a contained scope.

\paragraph{Reading the dependency map.} An institution that has
mature observability and identity but immature network segmentation
can reasonably proceed with Wave~4 (Policy) before completing
Wave~3 (Network), provided the policy enforcement boundary remains
within the segments already secured. The dependency map is a
constraint, not a strict ordering.

%-----------------------------------------------------------------------------
\section{Trajectory selection matrix}
\label{sec:pr-trajectories}

The trajectory through the waves is a function of the institution's
starting position, captured in the self-assessment of
\trchap{ch:self-assessment}{}. The matrix in
Table~\ref{tab:pr-trajectory-detail} expands the heuristics of
\S\ref{sec:sa-results} with their operational implications.

\begin{table}[H]
\centering\small
\caption{Trajectories through the waves, with their dominant
characteristics. The Selective and Deferred profiles are tactically
divergent and are treated as variations of the three principal
trajectories.}
\label{tab:pr-trajectory-detail}
\begin{tabularx}{\textwidth}{%
  >{\hsize=0.4\hsize\linewidth=\hsize\bfseries\raggedright\arraybackslash}X%
  >{\hsize=1.6\hsize\linewidth=\hsize\raggedright\arraybackslash}X%
}
\toprule
Trajectory & Dominant characteristics through the waves \\
\midrule
Greenfield                  & Rare; applies to new institutions or to a deliberate refoundation. The waves are sequential rather than coexistent because there is little to coexist with. \\
Brownfield single-suite-centric & A common starting point, especially for institutions whose estate is anchored on a single dominant productivity-suite incumbent. The playbook adds an open policy plane and an open observability collector layer in front of the existing incumbent-suite estate; that estate evolves at its own renewal pace. \\
Brownfield hybrid           & Frequent where heterogeneous estates and at least one open-source identity component already coexist. The playbook substitutes proprietary components at renewal moments while preserving the open-source perimeter; the institution is rarely starting from zero on any wave. \\
Selective adoption          & For institutions already largely aligned with the architecture. The playbook is used as a checklist, with waves run only on the non-aligned components. \\
Deferred                    & For institutions whose self-assessment shows heterogeneity without a dominant trend. Wave~0 is the first step, and the trajectory is re-baselined after that wave. \\
\bottomrule
\end{tabularx}
\end{table}

%-----------------------------------------------------------------------------
\section{Minimum viable subset for resource-constrained institutions}
\label{sec:pr-mvs}

The trajectories above describe \emph{how} an institution moves
through the full wave programme, and the pace options
(\S\ref{sec:pr-pace}) describe \emph{how fast}. A resource-constrained
institution may also legitimately adopt \emph{less}: a minimal
coherent core that delivers the size-invariant, legally compelled
functions of \trchap{ch:architectural-principles}{}
(\S\ref{sec:ap-necessity}) without committing to the full programme,
and that can be extended later along the same dependency graph.

\paragraph{The minimal core.} The dependency map
(Table~\ref{tab:pr-deps}) defines this subset precisely: Wave~0
(\trchap{ch:wave0-foundations}{} --- governance and inventory),
Wave~1 (\trchap{ch:wave1-identity}{} --- sovereign identity) and
Wave~2 (\trchap{ch:wave2-observability}{} --- observability and the
audit trail), with Wave~4 (\trchap{ch:wave4-policy}{} --- policy) a
natural next step once Identity and Observability are in place.
Waves~1 and~2 depend only on Wave~0, so the subset is internally
consistent and needs no later wave to function.

\paragraph{Why these waves.} This core delivers exactly the functions
that no institution offering a digital infrastructure can avoid:
documented governance and inventory, sovereign authentication, and the
observability and audit trail required under data-protection
obligations and, where the institution is in scope, NIS2. It is the
operational expression, in playbook terms, of the minimal necessary
set of \S\ref{sec:ap-necessity}.

\paragraph{What it does not deliver --- stated plainly.} The minimal
core does not provide the deeper network segmentation that contains a
breach (Wave~3, \trchap{ch:wave3-network}{}), managed endpoint fleet
control (Wave~6) or sovereign workload and storage (Wave~8,
\trchap{ch:wave8-workload}{}). A constrained institution defers these
or consumes them as mutualised services (\trchap{ch:sizing}{},
\S\ref{sec:sz-downscale}). The minimal subset is therefore a
legitimate \emph{steady state}, not an abandoned migration: the
conformance index (\trchap{ch:conformance}{}) is read as a maturity
trajectory, and the institution extends the subset wave by wave when
capacity allows.

%-----------------------------------------------------------------------------
\section{Pace options}
\label{sec:pr-pace}

The reference pace options describe three indicative horizons for
completing the playbook. They are operational planning aids, not
contractual commitments: the actual pace is set by the institution's
operational rhythm and by the renewal cycles of the components the
playbook substitutes.

\begin{description}[leftmargin=1.5em,style=nextline,nosep]
  \item[5-year compressed.] The playbook is run with significant
    parallelisation, with waves overlapping, and with substantial
    additional resourcing during the active period. This pace is
    appropriate for institutions whose strategic context warrants
    accelerated transformation (sovereignty risk crystallised by an
    incident, regulatory deadline, decisive executive mandate). It
    requires either substantial budget supplements or a willingness
    to defer other strategic initiatives for the duration. Risk
    tolerance is structurally higher.
  \item[7-year balanced.] The recommended default. Waves proceed
    largely in sequence with limited parallelisation, additional
    resourcing is moderate, and most waves close within their
    reversibility window without difficulty. This pace tolerates
    institutional rhythm changes (academic-year cadence, budget
    cycles, leadership transitions) without losing direction.
  \item[10-year gradual.] The playbook is run with minimal
    additional resourcing, with each wave aligned to the natural
    renewal cycle of the components it substitutes. Risk tolerance
    is structurally lower; coexistence periods are longer; the
    institution accumulates less migration debt at any moment.
    This pace is appropriate for institutions whose context does
    not warrant acceleration and whose financial planning operates
    on multi-year horizons.
\end{description}

\paragraph{Pace and trajectory interact.} A greenfield trajectory
on a 10-year horizon is unusual; a brownfield-hybrid on a 5-year
horizon is demanding. The combinations that an institution can
sustain are bounded by its self-assessment results in
\trchap{ch:self-assessment}{} (Section~E intent), and revising the
pace is a normal part of operating the playbook.

%-----------------------------------------------------------------------------
\section{Risk-and-reversibility philosophy}
\label{sec:pr-philosophy}

The playbook treats reversibility as a non-negotiable property of
every wave. The reasoning is structural: an institution that
commits to a transformation it cannot undo at acceptable cost has
ceased to govern the transformation. Three practical commitments
follow from the principle.

\begin{itemize}[leftmargin=1.5em,nosep]
  \item Every wave declares a \emph{reversibility window}: the
    period during which the wave can be reversed at acceptable
    cost. Outside this window, reversal is still possible but is
    treated as a new wave with its own planning.
  \item The \emph{coexistence pattern} of a wave is the principal
    mechanism of reversibility: as long as the incumbent component
    remains operationally viable, the wave can be reversed by
    re-routing to the incumbent.
  \item The \emph{exit criteria} of a wave include passage of the
    relevant conformance subset (\trchap{ch:conformance}{}); a
    wave that closes without the conformance subset is not
    closed in this playbook's sense, regardless of stakeholder
    declarations.
\end{itemize}

\paragraph{When reversibility is forfeited.} The principle is
non-negotiable in normal operation; it admits exceptions when the
institution's competent governing authority, at whatever level
holds the mandate to bind the institution, accepts the loss of
reversibility explicitly, in writing, against the sovereignty grid. The
architectural posture is that such forfeits should be exceptional
and that they should be revisited at every renewal as part of the
sovereignty grid's annual review.

\medskip

\noindent The remainder of Part~III develops each wave in turn.
Each wave chapter follows the same template: goal and dependencies,
pre-requisites, coexistence pattern, the three principal
trajectories, reversibility criteria, exit criteria, and a
real-world reference where one is publicly documented.
